\pagebreak

%
\begin{center}
\textbf{\large Supplementary}
\end{center}

\horizontalcherriesthree
\horizontalcherriesgrids
%

\section{Background}
\label{sec:background}

\paragraph{Diffusion Models}\citep{SohlDickstein2015DeepUL,Song2020ScoreBasedGM,ho2020denoising} generate data by approximating the reverse ODE to a stochastic forward process which transforms data to noise. They have become the standard approach for generative modeling of  images~\cite{dhariwal2021diffusion,ramesh2022hierarchical,saharia2022photorealistic,Rombach_2022,balaji2022ediffi} and videos~\citep{singer2022makeavideo,ho2022imagen,esser2023structure, blattmann2023align,gupta2023photorealistic}. Since these models can be derived both via a variational lower bound on the negative likelihood~\citep{SohlDickstein2015DeepUL} and score matching~\citep{Hyvrinen2005EstimationON,Vincent2011ACB,song2020generative}, 
%
various formulations of forward- and reverse processes~\citep{Song2020ScoreBasedGM,dockhorn2021score}, model parameterizations~\citep{ho2020denoising,ho2022classifierfree,Karras2022ElucidatingTD}, loss weightings~\citep{ho2020denoising, Karras2022ElucidatingTD} and ODE 
solvers~\citep{song2022denoising,lu2023dpmsolver, dockhorn2022genie} have led to a large number of different training objectives and sampling procedures. More recently, the seminal works of \citet{Kingma2023UnderstandingDO} and \citet{Karras2022ElucidatingTD} have proposed unified formulations and introduced new theoretical and practical insights for training~\cite{Karras2022ElucidatingTD,Kingma2023UnderstandingDO} and inference~\cite{Karras2022ElucidatingTD}. However, despite these improvements, the trajectories of common ODEs involve partly significant amounts of curvature~\citep{Karras2022ElucidatingTD,liu2022flow}, which requires increased amounts of solver steps and, thus, renders fast inference difficult. To overcome this, we adopt rectified flow models whose formulation allows for learning straight ODE trajectories.  

\paragraph{Rectified Flow Models}
\cite{liu2022flow,albergo2022building,lipman2023flow} approach generative modeling by
constructing a transport map between two distributions through an ordinary
differential equation (ODE). This approach has close connections to continuous
normalizing flows (CNF) \cite{Chen2018NeuralOD} as well as diffusion models.
Compared to CNFs, Rectified Flows and Stochastic Interpolants have the
advantage that they do not require simulation of the ODE during training.
Compared to diffusion models, they can result in ODEs that are faster to
simulate than the probability flow ODE \cite{Song2020ScoreBasedGM} associated
with diffusion models. Nevertheless, they do not result in optimal transport
solutions, and multiple works aim to minimize the trajectory curvature further
\cite{lee2023minimizing,alex2023improving,aramalex2023multisample}.
\cite{dao2023flow,ma2024sit} demonstrate the feasibility of rectified flow
formulations for class-conditional image synthesis, \cite{fischer2023boosting} for
latent-space upsampling, and \cite{liu2023instaflow}
apply the reflow procedure of \cite{liu2022flow} to distill a pretrained
text-to-image model \cite{Rombach_2022}. Here, we are interested in rectified
flows as the foundation for text-to-image synthesis with fewer sampling steps.
We perform an extensive comparison between different formulations and loss
weightings and propose a new timestep schedule for training of
rectified flows with improved performance.



\paragraph{Scaling Diffusion Models}
The transformer architecture~ \cite{vaswani2017attention} is well known for its scaling properties in NLP~\cite{kaplan2020scaling} and computer vision tasks~\cite{dosovitskiy2020image,zhai2022scaling}. For diffusion models, U-Net architectures~\cite{Ronneberger_2015} have been the  dominant choice~\cite{ho2020denoising,Rombach_2022,balaji2022ediffi}.
While some recent works explore diffusion transformer backbones~\cite{Peebles_2023,chen2023pixart,ma2024sit}, scaling laws for text-to-image diffusion models remain unexplored.

%

%
%
%
%
%
%
%
%
%

%
%
%
%
%
%



%
%
%
%
%
%
%
%
%
%
%
%

%
%
%
%
%
%


%
%
%
%
%
%

%
%
%
%
%
%
%
%
%
%
%
%
%
%
%
%


%
%
%
%
%
%
%
%
%
%
%
%
%
%
%
%
%
%
%

%
%
%
%
%
%
%
%
%

%
%
%
%
%
%
%
%
%
%
%
%
%
%
%
%
%
%
%
%
%
%
%
%
%
%
%
%
%
%
%
%
%
%
%
%
%
%
%
%

%
%
%
%
%
%
%
%
%
%
%

%
%
%
%
%
%
%
%
%
%
%
%

%
%
%
%
%
%
%
%
%
%
%
%
%
%
%
%
%
%
%
%
%
%
%
%
%
%
%
%
%
%
%
%
%
%
%
%
%
%

%
%
%
%
%
%
%
%
%
%
%
%
%
%
%
  
%
%
%
%
%
%
%
%
%
%
%
%
%
%
%
%
%
%
%
%
%
%
%

%
%

%
%
%
%
%
%
%
%
%
%
%
%
%
%
%

%
%
%
%
%
%
%
%
%
%
%
%
%

%
%
%
%
%
%
%

%
%
%
%
%
%
%
%

%
%
%
%

%
%
%
%
%
%
%
%
%
%
%

%
%
%
%
%
%
%
%
%
%
%
%
%
%
%
%
%
%

%
%

%
%
%
%

%
%
%
%
%
%
%
%
%
%
%

%
%
%
%
%
%
%
%


%
%
%
%
%
%
%
%
%
%
%
%
%
%
%
%
%
%
%

%
%
%
%
%
%
%
%
%
%
%
%
%
%
%
%
%
%
%

%
%
%
%
%
%
%
%
%
%
%

%
%
%
%
%
%
%
%
%
%
%
%
%

%
%

%
%
%
%
%
%
%
%
%
%
%
%
%
%
%
%
%
%
%
%
%
%
%
%
%

%
%
%
%
%
%
%

%
%
%
%
%
%
%
%
%
%
%
%
%
%
%
%
%
%
%
%
%
%
%
%
%
%
%
%
%
%
%






%

%
%
%
%

%
%
%
%
%
%
%
%
%
%
%

%
%
%
%
%
%
%
%
%
%
%
%
%
%
%
%
%
%

%
%
%
%
%
%
%
%
%
%
%
%
%
%
%
%
%
%
%
%
%

%
%
%
%
%
%
%
%
%
%
%

%
%
%
%
%
%
%
%
%
%
%
%
%

%
%
%

%

%
%
%
%
%
%
%

%

%
%
%
%

%
%
%
%

%
%
%
%
%
%
%

%
%
%
%
%
%
%
%
%
%
%
%
%
%
%
%
%
%

%
%
%
%
%
%
%
%
%
%
%
%
%
%
%
%
%

%
%
%
%
%
%
%
%
%
%
%
%

%
%
%
%

%
%
%
%
%
%
%
%
%
%
%

%
%
%
%
%
%
%
%


%
%
%
%
%
%
%
%
%
%
%
%
%
%
%
%
%
%
%

%
%
%
%
%
%
%
%
%
%
%
%


%
%
%
%
%
%
%
%
%
%
%



%
%
%
%
%
%
%
%
%
%
%
%
%
%
%
%
%
%
%
%
%
%
%
%
%
%
  %
  %
  %
  %
  %
  %
%
\newcommand*\diff{\mathop{}\!\mathrm{d}}
\newcommand*\Diff[1]{\mathop{}\!\mathrm{d^#1}}
\DeclarePairedDelimiter\bra{\langle}{\rvert}
\DeclarePairedDelimiter\ket{\lvert}{\rangle}
\DeclarePairedDelimiterX\braket[2]{\langle}{\rangle}{#1\,\delimsize\vert\,\mathopen{}#2}

\newpage
 \section{On Flow Matching}
 \subsection{Details on Simulation-Free Training of Flows}
 \label{subsec:flowproofs}

Following~\citep{lipman2023flow}, to see that $u_t(z)$ generates $p_t$, we note that the continuity equation provides a necessary and sufficient condition \cite{Villani2008OptimalTO}:
\begin{align}
     & \frac{d}{dt} p_t(x) + \nabla \cdot [p_t(x) v_t(x)] = 0 \leftrightarrow  \text{$v_t$ generates probability density path $p_t$}. \label{eq:cont_eq}
\end{align}

Therefore it suffices to show that
\begin{align}
    - \nabla \cdot [u_t(z) p_t(z)]
      &=
    - \nabla \cdot [\mathbb{E}_{\epsilon \sim
   \mathcal{N}(0,I)} u_t(z \vert \epsilon) \frac{p_t(z \vert
   \epsilon)}{p_t(z)}  p_t(z) ] \label{eq:marg_cont_line1}\\
     &=  \mathbb{E}_{\epsilon \sim \mathcal{N}(0,I)} - \nabla \cdot  [u_t(z \vert \epsilon) p_t(z \vert \epsilon) ] \label{eq:marg_cont_line2} \\
     &= \mathbb{E}_{\epsilon \sim \mathcal{N}(0,I)}  \frac{d}{dt} p_t(z \vert \epsilon) \label{eq:marg_cont_line3}  = \frac{d}{dt} p_t(z),
 \end{align}
 where we used the continuity equation \Cref{eq:cont_eq} for $u_t(z \vert \epsilon)$ in line \Cref{eq:marg_cont_line2} to \Cref{eq:marg_cont_line3} since $u_t(z \vert \epsilon)$ generates $p_t(z \vert \epsilon)$ and the definition of \Cref{eq:marginal_u} in line \Cref{eq:marg_cont_line1}

 The equivalence of objectives $\mathcal{L}_{FM} \leftrightharpoons \mathcal{L}_{CFM}$~\citep{lipman2023flow} follows from
 \begin{align}
 \mathcal{L}_{FM}(\Theta) &= \mathbb{E}_{t, p_t(z)}|| v_{\Theta}(z, t) - u_t(z)  ||_2^2 \\
     &=  \mathbb{E}_{t, p_t(z)}|| v_{\Theta}(z, t) ||_2^2  - 2 \mathbb{E}_{t, p_t(z)} \braket{v_{\Theta}(z, t) }{  u_t(z) } + c
     \label{eq:equiv_line2} \\
       &=\mathbb{E}_{t, p_t(z)}|| v_{\Theta}(z, t) ||_2^2  - 2 \mathbb{E}_{t, p_t(z | \epsilon),  p(\epsilon) } \braket{v_{\Theta}(z, t) }{u_t(z | \epsilon )} + c \label{eq:equiv_line3} \\
      &=   \mathbb{E}_{t, p_t(z | \epsilon), p(\epsilon) }|| v_{\Theta}(z, t) - u_t(z | \epsilon)  ||_2^2 + c'  =   \mathcal{L}_{CFM}(\Theta) + c'
 \end{align}
 where $c, c'$ do not depend on $\Theta$ and line \Cref{eq:equiv_line2} to line \Cref{eq:equiv_line3} follows from:
 \begin{align}
 \mathbb{E}_{p_t(z | \epsilon),  p(\epsilon) } \braket{v_{\Theta}(z, t) }{u_t(z | \epsilon )}
 &=  \int \diff z  \int \diff \epsilon p_t(z | \epsilon) p(\epsilon)\braket{v_{\Theta}(z, t) }{u_t(z | \epsilon )} \\
 %
 &=  \int \diff z p_t(z) \braket{v_{\Theta}(z, t) }{ \int \diff \epsilon \frac{ p_t(z| \epsilon)}{p_t(z)} p(\epsilon) u_t(z | \epsilon )} \label{eq:equiv_intergral_line2} \\
%
 &=  \int \diff z p_t(z)  \braket{v_{\Theta}(z, t) }{  u_t(z) } =  \mathbb{E}_{p_t(z)} \braket{v_{\Theta}(z, t) }{  u_t(z) } \label{eq:equiv_intergral_line3} 
 \end{align}
 where we extended with $\frac{p_t(z)}{p_t(z)}$ in line \Cref{eq:equiv_intergral_line2} and used the definition of \Cref{eq:marginal_u} in line \Cref{eq:equiv_intergral_line2} to \Cref{eq:equiv_intergral_line3}.

\subsection{Details on Image and Text Representations}
\label{suppsec:detailsonreps}
\textbf{Latent Image Representation}
We follow LDM \cite{Rombach_2022} and use a pretrained autoencoder to represent RGB images $X
\in \RR^{H\times W\times 3}$ in a smaller latent space $x=E(X) \in \RR^{h\times
w \times d}$. We use a spatial downsampling factor of $8$, such that
$h=\frac{H}{8}$ and $w=\frac{W}{8}$, and experiment with different values for
$d$ in \Cref{sec:improvedae}. We always apply the forward process
from \eqref{eq:forwardprocess} in the latent space, and when sampling a %
representation $x$ via \eqref{eq:ode}, we decode it back into pixel space $X =
D(x)$ via the decoder $D$.
We follow \citet{Rombach_2022} and normalize the latents by their mean and 
standard deviation, which are globally computed over a subset
of the training data. 
\Cref{fig:aefidstudy} shows how generative model training for different $d$ evolves as a function of model capacity, 
as discussed in \Cref{sec:improvedae}.
\aefidstudy
%

\textbf{Text Representation}
Similar to the encoding of images to latent representations, we also follow
previous approaches \cite{saharia2022photorealistic, balaji2022ediffi} and encode the text conditioning $c$ using
pretrained, frozen text models. In particular, for all experiments, we use a
combination of CLIP \cite{radford2021learning} models and a encoder-decoder text model.
Specifically, we encode $c$ with the text encoders of both a CLIP L/14 model of
\citet{radford2021learning} as well as an OpenCLIP bigG/14 model of \citet{Cherti_2023}. We
concatenate the pooled outputs, of sizes $768$ and $1280$ respectively, to
obtain a vector conditioning $c_\text{vec} \in \RR^{2048}$. We also concatenate
the penultimate hidden representations channel-wise to a CLIP context conditioning
$c_\text{ctxt}^{\text{CLIP}} \in \RR^{77\times2048}$. Next, we encode $c$ also
to the final hidden representation,
$c_\text{ctxt}^{\text{T5}}\in\RR^{77\times4096}$, of the encoder of a T5-v1.1-XXL model
\cite{raffel2019exploring}. Finally, we zero-pad $c^\text{CLIP}_\text{ctxt}$ along the channel
axis to $4096$ dimensions to match the T5 representation and concatenate it
along the sequence axis with $c_\text{ctxt}^\text{T5}$ to obtain the final
context representation $c_\text{ctxt}\in\RR^{154\times4096}$. These two caption
representations, $c_\text{vec}$ and $c_\text{ctxt}$, are used in two different
ways as described in \Cref{subsec:methodarch}.

\subsection{Preliminaries for the Experiments in \Cref{subsec:improvingrfs}.}
\label{subsubsec:exp_prelim}
%
%
%
%
%
%
%
%
%
%
%

\textbf{Datasets}
We use two datasets to account for the missing of a standard text-to-image
benchmark. As a widely used dataset, we convert the ImageNet dataset \cite{Russakovsky2014ImageNetLS}
into a dataset suitable for text-to-image models by adding captions of the
form ``a photo of a \texttt{\textlangle{}class name\textrangle{}}'' to images, where \texttt{\textlangle{}class name\textrangle{}} is randomly
chosen from one of the provided names for the image's class label.
%
As a more realistic text-to-image dataset, we use the CC12M dataset \cite{Changpinyo2021Conceptual1P} for
training.

\textbf{Optimization}
In this experiment, we train all models using a global batch size of 1024 using the AdamW optimizer
\cite{Loshchilov2017FixingWD} with a learning rate of $10^{-4}$ and 1000 linear warmup
steps. We use mixed-precision training and keep a copy of the model weights
which gets updated every 100 training batches with an exponential moving
average (EMA) using a decay factor of $0.99$. For unconditional diffusion
guidance \cite{ho2022classifierfree}, we set the outputs of each of the three text encoders
independently to zero with a probability of $46.4\%$, such that we roughly
train an unconditional model in $10\%$ of all steps.

\textbf{Evaluation}
%
%
%
%
%
%
As described in \Cref{subsec:improvingrfs}, we use CLIP scores, 
FID and validation losses to evaluate our models 
regularly during training on the COCO-2014 validation split~\cite{Lin_2014}.

As the loss values differ widely in magnitude and variance for different
timesteps, we evaluate them in a stratified way on eight equally spaced values
in the time interval $[0, 1]$.

To analyze how different approaches behave under different sampler settings, we
produce 1000 samples for each of the samplers which differ in guidance scales
as well as number of sampling steps. We evaluate these samples with CLIP scores
using CLIP L/14 \cite{radford2021learning} and also compute FID between CLIP L/14 image features of these
samples and the images of the validation set. For sampling, we always use a
Euler discretization \cite{euler1768institutionum} of \eqref{eq:ode} and six
different settings: 50 steps with classifier-free-guidance scales 1.0, 2.5,
5.0, and 5, 10, 25 steps with classifier-free-guidance scale 5.0.

\subsection{Improving SNR Samplers for Rectified Flow Models}
As described in \Cref{sec:method}, we introduce novel densities $\pi(t)$ for the timesteps
that we use to train our rectified flow models. \Cref{fig:tsamplers} visualizes the distributions of the \textbf{logit-normal sampler} and the \textbf{mode sampler} introduced in \Cref{subsubsec:tailoredsnrforrf}. Notably, as we demonstrate in \Cref{subsec:improvingrfs}, the \textbf{logit-normal sampler} outperforms the classic uniform rectified flow formulation~\citep{liu2022flow} and established diffusion baselines such as EDM~\citep{Karras2022ElucidatingTD} and LDM-Linear~\citep{Rombach_2022}.
\tsamplers
%

%
%
%
%
%
%
%

%

\figqualitativescaleparti

\section{Direct Preference Optimization}
\label{supsec:dpo}
\figdpo
\figdpohumaneval
Direct Preference Optimization~(DPO)~\citep{rafailov2023direct} is a technique to finetune LLMs with preference data. Recently, this method has been adapted to preference finetuning of text-to-image diffusion models~\citep{wallace2023diffusion}. In this section, we verify that our model is also amenable to preference optimization. In particular, we apply the method introduced in~\citet{wallace2023diffusion} to our 2B and 8B parameter base model. Rather than finetuning the entire model, we introduce learnable Low-Rank Adaptation (LoRA) matrices (of rank 128) for all linear layers as is common practice. We finetune these new parameters for 4k and 2k iteration for the 2B and 8B base model, respectively. We then evaluate the resulting model in a human preference study using a subset of 128 captions from the Partiprompts set~\citep{yu2022scaling} (roughly three voter per prompt and comparison). \Cref{fig:dpo_human_eval} shows that our base models can be effectively tuned for human preference. \Cref{fig:dpo} shows samples of the respective base models and DPO-finetuned models.

\section{Finetuning for instruction-based image editing}

A common approach for training instruction based image editing and general image-to-image diffusion models is to concatenate the latents of the input image to the noised latents of the diffusion target along the channel dimension before feeding the input into a U-Net ~\citep{brooks2023instructpix2pix, sheynin2023emu, saharia2022palette, saharia2022image}. We follow the same approach, concatenating input and target along the channels before patching, and demonstrate  that the same method is applicable to our proposed architecture. We finetune the 2B parameter base model on a dataset consisting of image-to-image editing tasks similar to the distribution of the InstructPix2Pix dataset \citep{brooks2023instructpix2pix} as well as inpainting, segmentation, colorization, deblurring and controlnet tasks similar to Emu Edit and Palette ~\citep{sheynin2023emu, saharia2022palette}. As shown in Fig \ref{fig:textmanipulation} we observe that the resulting 2B Edit model has the capability to manipulate text in a given image, even though no text manipulation tasks were included in the training data. We were not able to reproduce similar results when training a SDXL-based ~\citep{podell2023sdxl} editing model on the same data. 

\textmanipulation


\section{Data Preprocessing for Large-Scale Text-to-Image Training}
\label{suppsec:datepreproc}
%
%
%
%
%
%
%
%
%
%
%
%
%
%
%
%
%
%
%
%
%
%
%
%
%

\subsection{Precomputing Image and Text Embeddings}
\label{suppsec:precomputingembeddings}
%
%
%
Our model uses the output of multiple pretrained, frozen networks as inputs (autoencoder latents and text encoder representations). %
Since these outputs are constant during training, we precompute them once for the entire dataset.
%
%
This comes with two main advantages: (i) The encoders do not need to be available on the GPU during training, lowering the required memory. (ii) The forward encoding pass is skipped during training, saving time and total needed compute after the first epoch, see \cref{tab:preencodingspeedups}.
%

\preencodingspeedups

This approach has two disadvantages: First, random augmentation for each sample every epoch is not possible and we use square-center cropping during precomputation of image latents. For finetuning our model at higher resolutions, we specify a number of aspect ratio buckets, and resize and crop to the closest bucket first and then precompute in that aspect ratio. 
Second, the dense output of the text encoders is particularly large, creating additional storage cost and longer loading times during training (\cf \cref{tab:preencodingspeedups}).
We save the embeddings of the language models in half precision, as we do not observe a deterioration in performance in practice.

%
%

%
%
%


\subsection{Preventing Image Memorization}\label{sec:dedup}
%
In the context of generative image models memorization of training samples can lead to a number of issues~\citep{somepalli2023diffusion, carlini2023extracting, somepalli2023understanding}. 
To avoid verbatim copies of images by our trained models, we carefully scan our training dataset for duplicated examples and remove them. 
%

\paragraph{Details on Deduplication}
In accordance with the methods outlined by ~\citet{carlini2023extracting} and ~\citet{somepalli2023diffusion}, we opt for SSCD \citep{pizzi2022self} as the backbone for the deduplication process. The SSCD algorithm is a state-of-the-art technique for detecting near-duplicate images at scale, and it generates high-quality image embeddings that can be used for clustering and other downstream tasks. We also decided to follow~\citet{dalle2mitigations} to decide on a number of clusters $N$. For our experiments, we use $N = 16,000$.

We utilize~\citet{autofaiss} for clustering. ~\citet{autofaiss} is a library that simplifies the process of using Faiss (Facebook AI Similarity Search) for large-scale clustering tasks. Specifically, leverage FAISS index factory\footnote{\url{https://github.com/facebookresearch/faiss/wiki/The-index-factory}} functionality to train a custom index with predefined number of centroids. This approach allows for efficient and accurate clustering of high-dimensional data, such as image embeddings.

\Algref{alg:cluster_deduplication} details our deduplication approach. We ran an experiment to see how much data is removed by different SSCD threshold as shown in~\Figref{fig:sscd_ablation}. Based on these results we selected four thresholds for the final run~\Figref{fig:sscd_final_result}.

\subsection{Assessing the Efficacy of our Deduplication Efforts}
~\citet{carlini2023extracting} devise a two-stage data extraction attack that generates images using standard approaches, and flags those that exceed certain membership inference scoring criteria. \citet{carlini2023extracting} bias their search towards duplicated training examples because these are orders of magnitude more likely to be memorized than non-duplicated examples ~\citep{somepalli2023diffusion, somepalli2023diffusion,lee2021deduplicating}. 

%
To assess how well our SSCD-based deduplication works, we follow \citet{carlini2023extracting} to extract memorized samples from small, specifically for this purpose trained models and compare them before and after deduplication. Two main step of the mentioned procedure include: 1) Generate many examples using the diffusion model in the standard sampling manner and with the known prompts. 2) Perform membership inference to separate the model’s novel generations from those generations which are memorized training examples. \Algref{alg:memorization_detection} shows the steps to find the memorized samples based on~\citet{carlini2023extracting}. Note that we run this techniques two times; one for SD-2.1 model with only exact dedup removal as baseline, and for a model with the SD2.1 architecture but trained on removed exact duplication and near-duplication using SSCD~\citep{pizzi2022self}. 

We select the 350,000 most-duplicated examples from the training dataset based on SSCD~\citep{pizzi2022self} with threshold of 0.5, and generate 500 candidate images for each text prompt to increase the likelihood of finding memorization. The intuition is that for diffusion models, with high probability $Gen(p;r_1) \approx _d Gen(p;r_2)$ for two different random initial seeds $r_1$,$r_2$. On the other hand, if $Gen(p;r_1)\approx _d Gen(p;r_2)$ under some distance measure d, it is likely that these generated samples are memorized examples. To compute the distance measure $d$ between two images, we use a modified Euclidean $l_2$ distance. In particular, we found that many generations were often spuriously similar according to  $l_2$ distance (e.g., they all had gray backgrounds). We therefore instead divide each image into 16 non-overlapping 128 × 128 tiles and measure the maximum of the $l_2$ distance between any pair of image tiles between the two images. 
%
\Cref{fig:memorization_sscd} shows the comparison between number of memorized samples, before and after using SSCD with the threshold of 0.5 to remove near-duplicated samples. \citet{carlini2023extracting} mark images within clique size of 10 as memorized samples. Here we also explore different sizes for cliques. For all clique thresholds, SSCD is able to significantly reduce the number of memorized samples. Specifically, when the clique size is 10, trained SD models on the deduplicated training samples cut
off at SSCD$=0.5$ show a $5\times$ reduction in potentially memorized examples.

%
%

%

%

%
%


\algclusterdeduplication
%
%
\figdupcombined
\algmemorizationdetection
\figmembaselinessd
%
%
